% ============================================================
% Memoria del proyecto GPT2_MOE
% Formato: artículo científico a doble columna (estilo AAAI),
% 4 páginas de contenido + referencias
% ============================================================
\documentclass[10pt,twocolumn,a4paper]{article}

% ---------- Idioma y codificación ----------
\usepackage[utf8]{inputenc}
\usepackage[T1]{fontenc}
\usepackage[spanish,es-tabla,es-nodecimaldot,es-noshorthands]{babel}

% ---------- Tipografía y página ----------
\usepackage{mathptmx}                 % Times (texto y matemáticas)
\usepackage[scaled=0.92]{helvet}
\usepackage{courier}
\usepackage{microtype}
\usepackage[a4paper,textwidth=7in,textheight=9.6in,top=0.75in,columnsep=0.33in,hcentering]{geometry}
\setlength{\parindent}{10pt}
\frenchspacing

% ---------- Secciones (estilo AAAI) ----------
\usepackage{titlesec}
\titleformat{\section}{\centering\normalfont\large\bfseries}{\thesection}{0.6em}{}
\titleformat{\subsection}{\normalfont\normalsize\bfseries}{\thesubsection}{0.6em}{}
\titleformat{\paragraph}[runin]{\normalfont\normalsize\bfseries}{}{0pt}{}
\titlespacing*{\section}{0pt}{2ex plus .5ex minus .2ex}{1ex plus .2ex}
\titlespacing*{\subsection}{0pt}{1.6ex plus .5ex minus .2ex}{0.8ex plus .2ex}
\titlespacing*{\paragraph}{0pt}{1ex plus .3ex minus .2ex}{0.8em}

% ---------- Matemáticas, tablas y figuras ----------
\usepackage{amsmath,amssymb}
\usepackage{graphicx}
\usepackage{booktabs}
\usepackage{multirow}
\usepackage{enumitem}
\usepackage{xcolor}
\usepackage[font=small,labelfont=bf,labelsep=colon]{caption}

% ---------- Referencias y enlaces (autor-año, como AAAI) ----------
\usepackage[authoryear,round]{natbib}
\setcitestyle{aysep={}}
\usepackage[hidelinks,hyperfootnotes=false]{hyperref}
\usepackage[spanish,noabbrev]{cleveref}

% ============================================================
% DATOS DEL DOCUMENTO
% ============================================================
\newcommand{\papertitle}{TÍTULO DEL TRABAJO}
\newcommand{\paperdate}{25/10/2026}

\pagestyle{plain}

% ============================================================
\begin{document}

\twocolumn[{%
  \centering
  {\fontsize{14.4}{17}\selectfont\bfseries \papertitle\par}
  \vspace{1.2em}
  {\fontsize{11}{13}\selectfont\bfseries
    Carlos Molera Canals, Pablo García Beltrá, Linxi Jiang, Manuel García Romo\par}
  \vspace{0.4em}
  {\normalsize Universidad de Alicante\par}
  \vspace{2em}
}]

% Nota al pie de la primera columna (como el copyright de AAAI)
{\renewcommand{\thefootnote}{}\footnotetext{\raggedright\footnotesize Memoria del proyecto
GPT-2 MoE. Universidad de Alicante, \paperdate.}}

% ------------------------------------------------------------
% RESUMEN
% ------------------------------------------------------------
\begin{center}
  \bfseries Resumen
\end{center}
\vspace{-0.6em}
{\small\leftskip=0.25in\rightskip=0.25in\noindent

\par}

% ------------------------------------------------------------
\section{Introducción}
\label{sec:intro}

\paragraph{Trabajo relacionado}

\paragraph{Contribuciones}

% ------------------------------------------------------------
\section{Tareas y datos}
\label{sec:data}

Para estudiar si los expertos se especializan es necesario un conjunto de
tareas claramente diferenciadas y con una única respuesta correcta. Por ello construimos un dataset sintético con nuestro generador que incluye las seis tareas. Cada tarea se genera de forma
independiente con una semilla fija (42), por lo que el conjunto es
reproducible, y se guarda en un CSV con dos columnas, entrada
(\texttt{prompt}) y la respuesta esperada (\texttt{y}).

Cada ejemplo se convierte en la secuencia
\begin{center}
  \small\texttt{tarea: prompt = y}
\end{center}
de modo que el nombre de la tarea queda explícito al inicio del contexto (opcional). La
secuencia se tokeniza con el tokenizador BPE de GPT-2. La pérdida solo se calcula sobre los tokens
de la respuesta y el de fin de secuencia, el \emph{prompt} se enmascara. Se
descartan los ejemplos duplicados y los que superan la longitud de contexto
(64 tokens), y las secuencias se rellenan por la derecha hasta la longitud
del ejemplo más largo de entrenamiento. Cada tarea se divide por separado en
entrenamiento, validación y test (80/10/10) con la misma semilla, de modo que
ninguna fila se comparte entre particiones y todas las tareas mantienen la
misma proporción en cada una. Cada fase de entrenamiento utiliza 100\,000
ejemplos.

Las seis tareas cubren capacidades distintas: suma y producto, transformación de toda la secuencia de entrada,
acceso a posiciones concretas y clasificación binaria. Además, difieren en la
longitud de la respuesta, desde un único token en la paridad hasta más de
veinte en la ordenación o el deletreo, y en el vocabulario, que usan tanto dígitos como letras y símbolos. Esta heterogeneidad es la condición para que
tenga sentido specializar expertos, si todas las tareas exigieran el mismo
tipo de cómputo, no habría nada que repartir entre ellos. Cada una tiene una única respuesta correcta, lo que permite evaluar con
\emph{exact match} sin ambigüedad.

\paragraph{Tareas}
\begin{itemize}[leftmargin=*,nosep]
  \item \textbf{Suma y resta} (\texttt{add\_sub}): dos enteros
    $a, b \in [1, 10^4]$ unidos por $+$ o $-$ (la respuesta puede ser negativa).
  \item \textbf{Multiplicación y división} (\texttt{mul\_div}): producto de
    dos enteros en $[1, 10^4]$ o división exacta.
  \item \textbf{Ordenación alfabética} (\texttt{alphabetic}): de 3 a 25
    letras mayúsculas y minúsculas separadas por espacios, que deben
    devolverse ordenadas sin distinguir mayúsculas.
  \item \textbf{Deletreo} (\texttt{spell}): una cadena de 3 a 25 caracteres
    (letras, dígitos y signos de puntuación) que debe reescribirse carácter a
    carácter, separada por espacios.
  \item \textbf{Paridad} (\texttt{parity}): un entero en $[0, 10^4]$; la
    respuesta es 0 si es par y 1 si es impar.
  \item \textbf{Primero y último} (\texttt{first\_last}): una cadena de 3 a
    25 letras; la respuesta es la concatenación de su primer y su último
    carácter.
\end{itemize}

% ------------------------------------------------------------
\section{Arquitectura}
\label{sec:arch}

\subsection{Modelo base}
\label{sec:base}

\subsection{Capa MoE}
\label{sec:moe}

\subsection{Pérdida de balanceo de carga}
\label{sec:loss}

% ------------------------------------------------------------
\section{Experimentos}
\label{sec:experiments}

\subsection{E1: modelo denso}
\label{sec:e1}

\subsection{E2: MoE con expertos preentrenados por tarea}
\label{sec:e2}

\subsection{E3: MoE con entrenamiento conjunto}
\label{sec:e3}

% ------------------------------------------------------------
\section{Resultados}
\label{sec:results}

% ------------------------------------------------------------
\section{Discusión}
\label{sec:discussion}

% ------------------------------------------------------------
\section{Conclusiones}
\label{sec:conclusions}

% ============================================================
% REFERENCIAS (citar con \citep{clave} o \citet{clave})
% Formato: \bibitem[Apellido et~al.(Año)]{clave}
% ============================================================
% Las referencias empiezan siempre en la página 5 (4 de contenido)
\clearpage
\makeatletter
\@whilenum\value{page}<5\do{\null\clearpage}
\makeatother
\renewcommand{\refname}{Referencias}
\renewcommand{\bibfont}{\small}
\setlength{\bibsep}{3pt plus 1pt}
\begin{thebibliography}{99}

\bibitem[Fedus, Zoph y Shazeer(2022)]{fedus2022switch}
Fedus, W.; Zoph, B.; y Shazeer, N. 2022.
\newblock Switch Transformers: Scaling to Trillion Parameter Models with Simple and Efficient Sparsity.
\newblock \emph{Journal of Machine Learning Research}, 23(120): 1--39.

\bibitem[Jiang et~al.(2024)]{jiang2024mixtral}
Jiang, A.~Q.; Sablayrolles, A.; Roux, A.; Mensch, A.; Savary, B.; et~al. 2024.
\newblock Mixtral of Experts.
\newblock \emph{arXiv preprint arXiv:2401.04088}.

\bibitem[Komatsuzaki et~al.(2023)]{komatsuzaki2023upcycling}
Komatsuzaki, A.; Puigcerver, J.; Lee-Thorp, J.; Riquelme, C.; Mustafa, B.; Ainslie, J.; Tay, Y.; Dehghani, M.; y Houlsby, N. 2023.
\newblock Sparse Upcycling: Training Mixture-of-Experts from Dense Checkpoints.
\newblock En \emph{International Conference on Learning Representations (ICLR)}.

\bibitem[Shazeer et~al.(2017)]{shazeer2017moe}
Shazeer, N.; Mirhoseini, A.; Maziarz, K.; Davis, A.; Le, Q.; Hinton, G.; y Dean, J. 2017.
\newblock Outrageously Large Neural Networks: The Sparsely-Gated Mixture-of-Experts Layer.
\newblock En \emph{International Conference on Learning Representations (ICLR)}.

\end{thebibliography}

\end{document}
